\documentclass[11pt]{amsart}

\usepackage[margin=1in]{geometry}
\usepackage[T1]{fontenc}
\usepackage{lmodern}
\usepackage{microtype}
\usepackage{amsmath,amssymb,mathtools}
\usepackage{array,booktabs}
\usepackage{xcolor}
\usepackage{xurl}
\usepackage[colorlinks=true,
            linkcolor=blue!55!black,
            citecolor=blue!55!black,
            urlcolor=blue!55!black]{hyperref}

\newtheorem{theorem}{Theorem}[section]
\newtheorem{proposition}[theorem]{Proposition}
\newtheorem{lemma}[theorem]{Lemma}
\newtheorem{corollary}[theorem]{Corollary}
\theoremstyle{definition}
\newtheorem{definition}[theorem]{Definition}
\theoremstyle{remark}
\newtheorem{remark}[theorem]{Remark}

\newcommand{\Z}{\mathbf Z}
\newcommand{\F}{\mathbf F}
\newcommand{\Ga}{\mathbb G_a}
\newcommand{\Gm}{\mathbb G_m}
\newcommand{\Spec}{\operatorname{Spec}}
\newcommand{\Soc}{\operatorname{Soc}}
\newcommand{\gr}{\operatorname{gr}}
\newcommand{\length}{\operatorname{length}}
\newcommand{\m}{\mathfrak m}
\newcommand{\ot}{\otimes}
\newcommand{\id}{\mathrm{id}}
\newcommand{\unit}{\mathrm e}
\newcommand{\Np}{N}

\title[A rank-\(p^2\) counterexample for every prime]
      {A Rank-\(p^2\) Finite Flat Group Scheme\\
       Not Killed by \(p^2\), for Every Prime \(p\)}
\author{}
\date{July 12, 2026}
\subjclass[2020]{14L15, 16T05}
\keywords{finite flat group scheme, power word, affine group,
Oort--Tate polynomial, truncated logarithm}

\begin{document}

\begin{abstract}
For every prime \(p\), we construct an Artin local ring \(R_p\) and a
finite locally free, necessarily noncommutative, group scheme \(G_p/R_p\)
of rank \(p^2\) for which the \(p^2\)-th power word is not the unit word.
The base is
\[
 R_p=\Z[a,b]/(a^{p+1},b^{2p-1},a^pb^{p-1}+p).
\]
Put \(w_i=p^{-1}\binom pi\) and
\[
 W_c(T)=\sum_{i=1}^{p-1}w_i c^{i-1}T^i.
\]
The coordinate algebra is
\[
 A_p=R_p[U,V]/(F,G),
\]
where
\[
 \begin{split}
 F&=U^p-a b^{p-1}W_a(U)+b^pW_b(V),\\
 G&=V^p-a^pW_b(V).
 \end{split}
\]
Writing \(L=1+aU\), \(M=1+bV\), and
\(\lambda=LM^{-1}\), the coproduct is
\[
 \Delta(U)=U\ot1+\lambda\ot U,
 \qquad
 \Delta(V)=V\ot\lambda^{-1}+1\ot V.
\]
The two monic equations give the basis
\(\{U^iV^j:0\leq i,j<p\}\).  A divided Oort--Tate identity and one
truncated-logarithm calculation prove that the equations are preserved
by the displayed product law.  The power words are
\[
 [p^2]^\#(U)=p b^{p-1}UV^{p-1}\ne0,
 \qquad [p^2]^\#(V)=0,
 \qquad [p^3]=\unit.
\]
Nonvanishing follows from
\[
 p b^{p-1}=-a^pb^{2p-2}\ne0.
\]
For \(p=2\), the formulas specialize exactly to the previously
constructed rank-four example over
\(\Z[a,b]/(a^3,b^3,a^2b+2)\).
\end{abstract}

\maketitle
\tableofcontents

\section{Statement of the construction}

The theorem of Deligne says that a \emph{commutative} finite locally
free group scheme is killed by its rank.  Thus the requested example
cannot have a commutative group law.  As usual, its coordinate algebra
will still be a commutative algebra; noncommutativity means that its
coproduct is not cocommutative.

Fix a prime \(p\).  For \(1\leq i\leq p-1\), set
\begin{equation}\label{eq:wi}
 w_i=\frac1p\binom pi\in\Z.
\end{equation}
Introduce the two convenient polynomials
\begin{equation}\label{eq:S-W}
 S_p(X)=\sum_{i=1}^{p-1}w_iX^i,
 \qquad
 W_c(T)=\sum_{i=1}^{p-1}w_i c^{i-1}T^i.
\end{equation}
They are related by
\begin{equation}\label{eq:cW=S}
 cW_c(T)=S_p(cT),
 \qquad
 (1+X)^p=1+X^p+pS_p(X).
\end{equation}

Define
\begin{equation}\label{eq:base}
 R_p=\Z[a,b]/(a^{p+1},b^{2p-1},a^pb^{p-1}+p)
\end{equation}
and the two monic polynomials
\begin{equation}\label{eq:F-G}
 \begin{split}
 F(U,V)&=U^p-a b^{p-1}W_a(U)+b^pW_b(V),\\
 G(V)&=V^p-a^pW_b(V).
 \end{split}
\end{equation}
Let
\begin{equation}\label{eq:A}
 A_p=R_p[U,V]/(F,G),
 \qquad L=1+aU,\qquad M=1+bV.
\end{equation}

The relations will imply that \(L\) and \(M\) are units.  Put
\begin{equation}\label{eq:lambda}
 \lambda=LM^{-1}.
\end{equation}
The point of using the \emph{exact} inverse in this formula is that it
makes the character calculation formal for every prime.  No
prime-dependent correction to \(\lambda\) is needed.

\begin{theorem}[The arbitrary-prime counterexample]\label{thm:main}
For every prime \(p\), the following assertions hold.
\begin{enumerate}
\item The ring \(R_p\) is an Artin local complete intersection of
characteristic \(p^2\), residue field \(\F_p\), and length
\((p+1)(2p-1)\).  Moreover,
\[
 p b^{p-1}=-a^pb^{2p-2}\ne0.
\]
\item The algebra \(A_p\) is free of rank \(p^2\) over \(R_p\), with
basis
\[
 \mathcal B=\{U^iV^j:0\leq i,j<p\}.
\]
\item There is a Hopf algebra structure on \(A_p\) with
\(\varepsilon(U)=\varepsilon(V)=0\) and
\begin{equation}\label{eq:Delta-main}
 \Delta(U)=U\ot1+\lambda\ot U,
 \qquad
 \Delta(V)=V\ot\lambda^{-1}+1\ot V.
\end{equation}
The element \(\lambda\) is group-like.
\item If \(G_p=\Spec A_p\), then
\begin{equation}\label{eq:power-main}
 [p^2]^\#(U)=p b^{p-1}UV^{p-1}\ne0,
 \qquad
 [p^2]^\#(V)=0,
 \qquad
 [p^3]=\unit.
\end{equation}
In particular, \(G_p\) has rank \(p^2\) and is not killed by \(p^2\).
\item The closed fiber is
\[
 (G_p)_{\F_p}\simeq\alpha_p\times\alpha_p,
\]
but \(G_p\) itself is noncommutative.
\end{enumerate}
\end{theorem}

The rest of the note proves the theorem.  The only substantial closure
calculation is isolated in Section~\ref{sec:closure}.  Sections
\ref{sec:base} and~\ref{sec:powers} give independent, transparent
proofs of the two facts that matter most: the coefficient in
\eqref{eq:power-main} is nonzero, and it is exactly the \(p^2\)-th power
carry.

\section{The Artin base and its top class}\label{sec:base}

Let
\[
 T=\Z[a,b]/(a^{p+1},b^{2p-1}).
\]
As an abelian group, \(T\) is free on the box monomials
\begin{equation}\label{eq:box}
 a^ib^j,\qquad 0\leq i\leq p,\quad 0\leq j\leq2p-2.
\end{equation}
Let \(N\) denote multiplication by \(a^pb^{p-1}\).  It has exactly the
following \(p\) nonzero arrows:
\begin{equation}\label{eq:arrows}
 b^j\longmapsto a^pb^{p-1+j},
 \qquad 0\leq j\leq p-1.
\end{equation}
All other box monomials are killed by \(N\), and \(N^2=0\).

The additive group of \(R_p\) is the cokernel of \(p\id+N\).  On each
two-term chain in \eqref{eq:arrows}, if \(e_j=b^j\) and
\(f_j=a^pb^{p-1+j}\), the relations are
\[
 pe_j+f_j=0,\qquad pf_j=0.
\]
Thus this chain contributes one copy of \(\Z/p^2\), generated by
\(e_j\).  Every remaining box monomial contributes one copy of
\(\Z/p\).  We have proved the precise Smith form
\begin{equation}\label{eq:additive}
 (R_p)_{\mathrm{add}}\simeq
 (\Z/p^2)^p\oplus
 (\Z/p)^{(p+1)(2p-1)-2p}.
\end{equation}

Several consequences are immediate.  First, \(R_p\) has characteristic
exactly \(p^2\).  Second,
\begin{equation}\label{eq:top-nonzero}
 p b^{p-1}=-a^pb^{2p-2}\ne0;
\end{equation}
it is the last element in the chain beginning with \(b^{p-1}\).  Third,
all of \(a,b,p\) are nilpotent, so \(R_p\) is local with maximal ideal
\(\m=(a,b)\) and residue field \(\F_p\).  Its length is
\[
 2p+\bigl((p+1)(2p-1)-2p\bigr)=(p+1)(2p-1).
\]

The associated graded ring has the especially simple form
\begin{equation}\label{eq:gr-base}
 \gr_{\m}R_p\simeq
 \F_p[A,B]/(A^{p+1},B^{2p-1}).
\end{equation}
In particular, \(a^pb^{2p-2}\) is the unique top monomial.  One can also
see directly that the three equations in \eqref{eq:base} form a regular
sequence.  Indeed, the first two do so, and multiplication by
\(p\id+N\) is injective on the torsion-free abelian group underlying
\(T\): if \((p\id+N)x=0\), applying \(p\id-N\) gives \(p^2x=0\), hence
\(x=0\).  Thus \(R_p\) is an Artin local complete intersection.

We record three scalar consequences used repeatedly below:
\begin{equation}\label{eq:scalar-zeroes}
 p^2=0,\qquad pa=0,\qquad pb^p=0.
\end{equation}
They follow respectively by multiplying
\(p=-a^pb^{p-1}\) by \(p,a,b^p\).

\section{Freeness and the two character identities}

First quotient \(R_p[V]\) by the monic polynomial \(G(V)\).  The result
is free over \(R_p\) on \(1,V,\ldots,V^{p-1}\).  Over that quotient,
the polynomial \(F(U,V)\) is monic of degree \(p\) in \(U\).  Therefore
\begin{equation}\label{eq:basis}
 A_p=\bigoplus_{0\leq i,j<p}R_pU^iV^j.
\end{equation}
This proves finite local freeness and rank \(p^2\) without any
Gr\"obner-basis calculation.

Put
\[
 x=bV,\qquad t=aU.
\]
The second relation in \eqref{eq:F-G} gives
\begin{align*}
 x^p
 &=b^pV^p
 =a^pb^pW_b(V)
 =a^pb^{p-1}S_p(x)
 =-pS_p(x).
\end{align*}
Consequently, by \eqref{eq:cW=S},
\begin{equation}\label{eq:Mp}
 M^p=(1+x)^p=1+x^p+pS_p(x)=1.
\end{equation}

Similarly, multiplying \(F=0\) by \(a^p\) gives
\begin{equation}\label{eq:tp}
 t^p=pS_p(x).
\end{equation}
Since \(pa=0\), every term of \(pS_p(t)\) vanishes.  Hence
\begin{equation}\label{eq:Lp}
 L^p=(1+t)^p=1+pS_p(x).
\end{equation}

Equations \eqref{eq:Mp} and \eqref{eq:Lp} show explicitly that \(M\)
and \(L\) are units.  In fact,
\begin{equation}\label{eq:explicit-inverses}
 M^{-1}=M^{p-1},
 \qquad
 L^{-1}=L^{p-1}\bigl(1-pS_p(x)\bigr),
\end{equation}
because \((pS_p(x))^2=0\).  Thus \(\lambda=LM^{-1}\) and
\(\lambda^{-1}=ML^{-1}\) are honest elements of \(A_p\), with no
localization implicit in the definition.

We will also use
\begin{equation}\label{eq:lambdap}
 \lambda^p=1+pS_p(x),
 \qquad
 \lambda^{p^2}=1.
\end{equation}
The second equality follows from the first and \(p^2=0\).

\section{The rank-\texorpdfstring{\(p\)}{p} affine identity}

The degree-\(p\) polynomials in the construction are the natural
mixed-characteristic replacements for the quadratic equations in the
rank-four example.

\begin{lemma}[Divided affine, or Oort--Tate, identity]
\label{lem:OT}
Let \(C\) be a ring, let \(c,d\in C\) satisfy
\(c^{p-1}d=p\), and put
\[
 \Phi_{c,d}(T)=T^p+dW_c(T).
\]
For
\[
 X\oplus_cY=X+Y+cXY
\]
one has the exact polynomial identity
\begin{equation}\label{eq:OT}
 \Phi_{c,d}(X\oplus_cY)
 =\Phi_{c,d}(X)+\Phi_{c,d}(Y)
  +c^p\Phi_{c,d}(X)\Phi_{c,d}(Y).
\end{equation}
\end{lemma}

\begin{proof}
The defining choice of coefficients gives
\begin{equation}\label{eq:OT-character}
 (1+cT)^p=1+c^p\Phi_{c,d}(T).
\end{equation}
On the other hand,
\[
 1+c(X\oplus_cY)=(1+cX)(1+cY).
\]
Raise this equality to the \(p\)-th power and apply
\eqref{eq:OT-character} to the three occurrences of \(\Phi\).  After
expansion one obtains \eqref{eq:OT}.  This is an integral polynomial
identity; no cancellation by a possibly nonregular element is being
used in the quotient algebra.
\end{proof}

For our first equation, take
\[
 c=a,\qquad d=-ab^{p-1};
\]
for the second, take
\[
 c=b,\qquad d=-a^p.
\]
In both cases \(c^{p-1}d=p\) in \(R_p\).  Thus
\begin{equation}\label{eq:Phiab}
 \Phi_a(T)=T^p-ab^{p-1}W_a(T),
 \qquad
 \Phi_b(T)=T^p-a^pW_b(T),
\end{equation}
and our relations are
\begin{equation}\label{eq:F-Phi}
 F=\Phi_a(U)+b^{p-1}S_p(bV),
 \qquad G=\Phi_b(V).
\end{equation}

\section{Closure under the affine-group law}\label{sec:closure}

We now prove that \eqref{eq:Delta-main} respects the two monic
relations.  This is the step at which it would be invalid merely to
multiply a relation defect by \(a^p\) or \(b^p\), since those elements
are zero divisors.  The argument below proves the actual monic defects
to be zero.

In \(A_p\ot_{R_p}A_p\), use subscripts \(1,2\) for the two factors and
put
\begin{equation}\label{eq:star}
 U_*=U_1+\lambda_1U_2,
 \qquad
 V_*=V_1\lambda_2^{-1}+V_2.
\end{equation}
We must prove
\begin{equation}\label{eq:three-closure}
 F(U_*,V_*)=0,
 \qquad G(V_*)=0,
 \qquad \lambda(U_*,V_*)=\lambda_1\lambda_2.
\end{equation}

\subsection{The second monic relation}

Since \(\lambda_2^{-1}=M_2L_2^{-1}\), one may write
\begin{equation}\label{eq:Vstar-formal}
 V_*=\bigl(L_2^{-1}V_1\bigr)\oplus_bV_2.
\end{equation}
Lemma~\ref{lem:OT} reduces the problem to showing
\(\Phi_b(L_2^{-1}V_1)=0\).  Using \(\Phi_b(V_1)=0\), a direct scaling
calculation gives
\begin{align}
 \Phi_b(L_2^{-1}V_1)
 &=a^p\sum_{i=1}^{p-1}w_i b^{i-1}V_1^i
   \bigl(L_2^{-p}-L_2^{-i}\bigr).\label{eq:scaled-G}
\end{align}
For every \(i\), the parenthesis is divisible by
\(L_2-1=aU_2\).  Its product with \(a^p\) is therefore zero because
\(a^{p+1}=0\).  Hence \eqref{eq:scaled-G} vanishes.  The product term
in \eqref{eq:OT} also vanishes, and we conclude that
\begin{equation}\label{eq:G-closes}
 G(V_*)=0.
\end{equation}

\subsection{The bridge and the truncated logarithm}

The first relation contains the bridge
\(b^{p-1}S_p(bV)\).  We isolate the elementary identity which makes it
close.

\begin{lemma}[Divided bridge lemma]\label{lem:bridge}
Assume \(p\) is odd.  In the tensor square above, set
\[
 x=bV_1,\qquad y=bV_2,\qquad t=aU_2,
 \qquad q=(1+x)^{-1},
\]
and
\begin{equation}\label{eq:z}
 z=bV_*=y+\frac{x(1+y)}{1+t}.
\end{equation}
Then
\begin{equation}\label{eq:E}
 \begin{split}
 E={}&S_p(z)-S_p(x)-S_p(y)
       +S_p(t)-S_p(tq)\\
    &\hspace{25mm}-pS_p(x)S_p(y)
 \end{split}
\end{equation}
satisfies
\begin{equation}\label{eq:bridge-vanish}
 b^{p-1}E=0.
\end{equation}
\end{lemma}

\begin{proof}
Modulo \(p\), the coefficients in \eqref{eq:wi} satisfy
\[
 w_i\equiv\frac{(-1)^{i-1}}{i}\pmod p.
\]
Thus \(S_p(T)\bmod p\) is the logarithm
\[
 \log(1+T)=\sum_{i=1}^{p-1}
             \frac{(-1)^{i-1}}{i}T^i
\]
truncated just before degree \(p\).  The two elementary factorizations
\begin{equation}\label{eq:cross-ratio}
 1+z=\frac{(1+y)(1+x+t)}{1+t},
 \qquad
 1+tq=\frac{1+x+t}{1+x}
\end{equation}
give the exact formal-logarithm identity
\begin{equation}\label{eq:log-identity}
 \begin{split}
 0={}&\log(1+z)-\log(1+x)-\log(1+y)\\
    &+\log(1+t)-\log(1+tq).
 \end{split}
\end{equation}

Compare coefficients in \(x,y\) of total degree less than \(p\), using
the truncated series above.  Equation \eqref{eq:log-identity} shows
that all such coefficients in the first five terms of \(E\) vanish
modulo \(p\), apart from terms containing \(t^p\).  Equivalently,
before substituting the scaled defining relations,
\begin{equation}\label{eq:E-membership}
 E\in p(x,y)+(x,y)^p+t^p(x,y).
\end{equation}
The last term in \eqref{eq:E} is already in \(p(x,y)\).  From
\eqref{eq:tp}, applied in the second tensor factor,
\(t^p=pS_p(y)\), so the last summand in
\eqref{eq:E-membership} also lies in \(p(x,y)\).

Finally, every element of \(p(x,y)\) is divisible by \(pb\), so
\[
 b^{p-1}p(x,y)=0
\]
by \(pb^p=0\).  Every monomial in \((x,y)^p\) is divisible by \(b^p\),
so
\[
 b^{p-1}(x,y)^p=0
\]
by \(b^{2p-1}=0\).  This proves \eqref{eq:bridge-vanish}.
\end{proof}

\begin{proposition}[Preservation of the first relation]
\label{prop:F-closes}
For odd \(p\), one has \(F(U_*,V_*)=0\).
\end{proposition}

\begin{proof}
Write \(q=M_1^{-1}\).  Then
\[
 U_*=U_1\oplus_a(qU_2).
\]
The scaling defect of \(\Phi_a\) is
\begin{equation}\label{eq:scaled-F}
 \Phi_a(qU_2)
 =\Phi_a(U_2)+b^{p-1}\bigl(S_p(t)-S_p(tq)\bigr),
\end{equation}
because \(q^p=1\).  Apply Lemma~\ref{lem:OT}, use
\[
 \Phi_a(U_i)=-b^{p-1}S_p(bV_i),
\]
and add the new bridge \(b^{p-1}S_p(bV_*)\).  The result is the exact
defect formula
\begin{equation}\label{eq:F-defect}
 F(U_*,V_*)=b^{p-1}E,
\end{equation}
with \(E\) as in \eqref{eq:E}.  The product term in the Oort--Tate
identity accounts for the final term
\(-pS_p(x)S_p(y)\); terms containing the scaling correction in
\eqref{eq:scaled-F} acquire one additional factor of \(x\) and die by
\(b^{2p-1}=0\).  Lemma~\ref{lem:bridge} now proves that
\eqref{eq:F-defect} is zero.
\end{proof}

\subsection{The character is group-like}

The character calculation is shorter than either relation calculation.
Put
\begin{equation}\label{eq:C}
 C=M_1+L_2-1=1+bV_1+aU_2.
\end{equation}
This is a unit: its difference from \(1\) is nilpotent.  Directly from
\eqref{eq:star},
\begin{align}
 L_*:=1+aU_*
 &=L_1M_1^{-1}C=\lambda_1C,\label{eq:Lstar}\\
 M_*:=1+bV_*
 &=M_2L_2^{-1}C=\lambda_2^{-1}C.\label{eq:Mstar}
\end{align}
Therefore
\begin{equation}\label{eq:lambda-group-like}
 \lambda(U_*,V_*)=L_*M_*^{-1}
 =\lambda_1\lambda_2.
\end{equation}

Equations \eqref{eq:G-closes}, \eqref{eq:F-defect}, and
\eqref{eq:lambda-group-like} prove all three identities in
\eqref{eq:three-closure} for odd \(p\).  Thus \eqref{eq:Delta-main}
defines an algebra map \(A_p\to A_p\ot A_p\).

Coassociativity is now formal.  Indeed, \(\lambda\) is group-like, \(U\)
is left \(\lambda\)-skew primitive, and \(V\) is right
\(\lambda^{-1}\)-skew primitive.  The formulas
\(\varepsilon(U)=\varepsilon(V)=0\), \(\varepsilon(\lambda)=1\) give
the counit identities.  We have constructed a finite flat affine monoid.
Section~\ref{sec:powers} will show that every element has a functorial
inverse, completing the Hopf-algebra proof.

\section{The \texorpdfstring{\(p^2\)}{p-squared}-th power carry}
\label{sec:powers}

For \(n\geq1\), put
\begin{equation}\label{eq:norm}
 N_n(T)=1+T+\cdots+T^{n-1}.
\end{equation}
The skew-primitive formulas give, in any associative monoid,
\begin{equation}\label{eq:power-skew}
 [n]^\#(U)=N_n(\lambda)U,
 \qquad
 [n]^\#(V)=V N_n(\lambda^{-1}).
\end{equation}
This is simply the usual geometric-sum formula for powers of an affine
matrix.

Assume first that \(p\) is odd.  From \eqref{eq:lambdap},
\[
 N_{p^2}(\lambda)=N_p(\lambda)N_p(\lambda^p).
\]
Since \(\lambda^p=1+pS_p(x)\), \(p^2=0\), and \(p\) is odd,
\begin{equation}\label{eq:Nplambdap}
 N_p(\lambda^p)=p.
\end{equation}
Moreover, \(pa=0\) implies \(pL^j=p\), and hence
\begin{equation}\label{eq:p-lambda}
 p\lambda^j=pM^{-j}.
\end{equation}
Because \(M^p=1\), inversion permutes the powers
\(1,M,\ldots,M^{p-1}\).  Therefore
\begin{align}
 N_{p^2}(\lambda)
 &=pN_p(\lambda)
 =pN_p(M)\nonumber\\
 &=p(bV)^{p-1}.
 \label{eq:Np2}
\end{align}
In the last step we used
\[
 N_p(1+x)=x^{p-1}
 +p\sum_{i=1}^{p-1}w_i x^{i-1}
\]
and \(p^2=0\).

The same argument gives
\begin{equation}\label{eq:Np2-inverse}
 N_{p^2}(\lambda^{-1})=p(bV)^{p-1}.
\end{equation}
Substituting \eqref{eq:Np2} and \eqref{eq:Np2-inverse} into
\eqref{eq:power-skew} yields
\begin{align}
 [p^2]^\#(U)&=p b^{p-1}UV^{p-1},\label{eq:p2U}\\
 [p^2]^\#(V)&=p b^{p-1}V^p=0.\label{eq:p2V}
\end{align}
For the final equality, use \(V^p=a^pW_b(V)\) and \(pa=0\).

The first element in \eqref{eq:p2U} is nonzero.  Indeed,
\(UV^{p-1}\) belongs to the \(R_p\)-basis \eqref{eq:basis}, while its
coefficient \(pb^{p-1}\) is nonzero by \eqref{eq:top-nonzero}.  This is
the desired failure of the rank-annihilation identity.

On the other hand, \(\lambda^{p^2}=1\), and
\begin{equation}\label{eq:Np3}
 N_{p^3}(\lambda)
 =N_{p^2}(\lambda)N_p(\lambda^{p^2})
 =pN_{p^2}(\lambda)=0.
\end{equation}
The same holds for \(\lambda^{-1}\).  Hence
\begin{equation}\label{eq:p3zero}
 [p^3]^\#(U)=[p^3]^\#(V)=0,
 \qquad [p^3]=\unit.
\end{equation}

The \((p^3-1)\)-st power word is consequently a two-sided inverse in
the finite flat monoid constructed above.  It supplies the antipode and
proves that \(A_p\) is a Hopf algebra.

\section{The prime two and the previous rank-four example}

When \(p=2\), one has \(w_1=1\), \(W_c(T)=T\), and the formulas become
\begin{equation}\label{eq:p2-base-A}
 R_2=\Z[a,b]/(a^3,b^3,a^2b+2),
\end{equation}
\begin{equation}\label{eq:p2-A}
 A_2=R_2[U,V]/(U^2-abU+b^2V,\ V^2-a^2V).
\end{equation}
The relation \(M^2=1\) gives \(M^{-1}=M\), so
\begin{equation}\label{eq:p2-lambda}
 \lambda=LM=(1+aU)(1+bV).
\end{equation}
This is exactly the earlier human-scale rank-four construction.

For completeness, the five scalar cancellations behind closure are
\[
 (ab)a=-2,\qquad (a^2)b=-2,\qquad
 a^2a=0,\qquad (-b^2)b=0,
\]
and
\[
 (-b^2)a+(ab)b=0.
\]
They imply directly that
\[
 U_*=U_1+\lambda_1U_2,
 \qquad V_*=V_1\lambda_2+V_2
\]
preserve both quadrics and make \(\lambda\) multiplicative.

Our uniform formula uses \(\lambda^{-1}\) in the second skew direction.
It gives the same coproduct here.  Indeed, if
\(\theta=\lambda-1\), then
\[
 \theta^2=0,\qquad 2\theta=2bV,
\]
so
\[
 V(\lambda^{-1}-\lambda)
 =-2V\theta=-2bV^2=-2a^2bV=4V=0.
\]

The power calculation is
\begin{equation}\label{eq:p2-powers}
 [4]^\#(U)=2bUV\ne0,
 \qquad [4]^\#(V)=0,
 \qquad [8]=\unit.
\end{equation}
Thus Theorem~\ref{thm:main} includes \(p=2\) literally, while the
truncated-log proof in Section~\ref{sec:closure} was needed only for odd
primes.

\section{Special fiber and noncommutativity}

Modulo the maximal ideal \(\m=(a,b)\), the two equations become
\[
 U^p=0,\qquad V^p=0,
\]
and \(\lambda=1\).  Both \(U\) and \(V\) become primitive.  Therefore
\begin{equation}\label{eq:special-fiber}
 (G_p)_{\F_p}\simeq\alpha_p\times\alpha_p.
\end{equation}

The total group is nevertheless noncommutative.  Since \(M^p=1\), one
may write
\[
 \lambda=LM^{p-1}.
\]
The coefficient of \(V\) in \(\lambda-1\) is
\((p-1)b\), which is nonzero in \(R_p\).  Now
\begin{equation}\label{eq:noncocomm}
 \Delta(U)-\tau\Delta(U)
 =(\lambda-1)\ot U-U\ot(\lambda-1).
\end{equation}
The coefficient of \(V\ot U\) in \eqref{eq:noncocomm} is
\((p-1)b\ne0\).  Since \(A_p\ot A_p\) is free on tensor products of
the basis \eqref{eq:basis}, the displayed element is nonzero.  Hence the
coproduct is not cocommutative and \(G_p\) is noncommutative, as it must
be in view of Deligne's theorem.

\section{The first odd case: \texorpdfstring{\(p=3\)}{p=3}}

It is useful to display the first genuinely new case without summation
notation.  Here
\[
 w_1=w_2=1,
 \qquad W_c(T)=T+cT^2.
\]
The base is
\begin{equation}\label{eq:p3-base}
 R_3=\Z[a,b]/(a^4,b^5,a^3b^2+3),
\end{equation}
and
\begin{equation}\label{eq:p3-A}
 \begin{split}
 A_3=R_3[U,V]/(&U^3-ab^2U-a^2b^2U^2+b^3V+b^4V^2,\\
               &V^3-a^3V-a^3bV^2).
 \end{split}
\end{equation}
It is free on the nine monomials \(U^iV^j\), \(0\leq i,j\leq2\).
With
\[
 L=1+aU,\qquad M=1+bV,
\]
the relations give
\[
 M^3=1,
 \qquad L^3=1+3bV+3b^2V^2.
\]
Thus
\[
 \lambda=LM^2,
 \qquad
 \lambda^{-1}=ML^2(1-3bV-3b^2V^2).
\]
The coproduct is
\[
 \Delta(U)=U\ot1+\lambda\ot U,
 \qquad
 \Delta(V)=V\ot\lambda^{-1}+1\ot V.
\]
The ninth-power calculation is
\begin{equation}\label{eq:p3-power}
 [9]^\#(U)=3b^2UV^2=-a^3b^4UV^2\ne0,
 \qquad [9]^\#(V)=0,
 \qquad [27]=\unit.
\end{equation}
The nonzero coefficient \(3b^2=-a^3b^4\) is the top class of
\(R_3\).

\section{Conceptual summary}

The construction has four moving parts.
\begin{enumerate}
\item The base relation
\[
 a^pb^{p-1}=-p
\]
turns the two degree-\(p\) polynomials into affine, Oort--Tate-type
blocks.  The deeper nilpotence bounds retain the final class
\[
 pb^{p-1}=-a^pb^{2p-2}.
\]
\item The term \(b^pW_b(V)\) in the first equation is the higher-prime
bridge.  Without it, the two rank-\(p\) directions would form an overly
simple extension.  The bridge is exactly what the divided logarithm in
Lemma~\ref{lem:bridge} transports under multiplication.
\item The ratio
\[
 \lambda=\frac{1+aU}{1+bV}
\]
is a character.  The opposite exact skew directions in
\eqref{eq:Delta-main} make its multiplicativity reduce to the common
factor \(C=M_1+L_2-1\).
\item The first \(p\)-adically nontrivial term in the geometric sum of
length \(p^2\) is
\[
 p(bV)^{p-1}.
\]
Multiplying by the skew coordinate \(U\) produces the surviving class
\(pb^{p-1}UV^{p-1}\).  One further block of \(p\) powers kills it, which
explains \([p^3]=\unit\).
\end{enumerate}

Thus the rank-four fourth-power carry is not an isolated phenomenon at
the prime two.  It is the first case of a uniform family of rank-\(p^2\)
counterexamples, with the square-zero expression \(2bUV\) replaced by
the divided \(p\)-stage carry
\[
 p b^{p-1}UV^{p-1}.
\]

\section{Independent exact checks}

The proof above is by hand.  As bounded local checks, the three closure
normal forms
\[
 F(U_*,V_*),\qquad G(V_*),\qquad
 \lambda(U_*,V_*)-\lambda_1\lambda_2
\]
were also reduced exactly over \(\Z\) in Macaulay2 for
\(p=3,5,7\).  In every case the three normal forms were
\[
 0,\qquad0,\qquad0.
\]
The \(p=5\) and \(p=7\) checks are useful guards against two accidental
small-prime simplifications: the coefficients \(w_i\) cease to be all
one, and the nilpotence box becomes substantially larger.  These checks
are not used as a substitute for the arbitrary-\(p\) divided identities
in Sections~\ref{sec:closure} and~\ref{sec:powers}.

The saved representative odd-prime audit is
\begin{center}
\path{m2/verify_rank_p2_uniform_p5_20260712.m2}.
\end{center}
It also verifies the nonzero \([25]\)-carry and the vanishing needed for
the \([125]\)-word.

For \(p=2\), the pre-existing exact audit is
\begin{center}
\path{m2/verify_human_scale_affine_rank4_counterexample_20260712.m2}.
\end{center}

\end{document}
